\title{SVD-Compressed Cross Sections in Monte Carlo
Neutron Transport\\[0.3em]
\normalsize An implementation-led benchmark, what
worked, what did not, and where compression actually pays}

\author{F.~Rodrigues\\
\small\texttt{sorcerer86pt@gmail.com}}
\date{\today}
\maketitle

\begin{abstract}
A pure-Rust Monte Carlo neutron transport engine implements
four cross-section providers behind a common interface ---
pointwise table, truncated SVD of $\log_{10}\sigma(E,T)$,
ACE$+$WMP, and an SVD$+$WMP hybrid --- and a matching CUDA
path for the pointwise and SVD providers. The question this
paper sets out to answer is whether truncated SVD is a
viable drop-in replacement for pointwise tables in
continuous-energy MC; the answer, on the measurements
collected in this engine, is \emph{only in a narrow regime,
and not the one the representation-byte literature points to}.

\paragraph{What we measured and what it actually does.}
At the operating point most production codes run --- a
single library temperature column per nuclide, full transport
physics including S($\alpha,\beta$) and URR --- the SVD
representation is monotonically \emph{larger} than the
pointwise table at every rank measured ($1$, $2$, $3$, $5$,
$7$). On the $9$-nuclide PWR pin cell at rank $5$, the
in-engine SVD working set is $527$~MB against $103$~MB for the
table ($5.12\times$ larger) and $100$~MB for ACE$+$WMP. CPU
throughput on the same problem favours the table by
$1.25\times$: the per-XS-lookup speedup that SVD shows in a
micro-benchmark ($8$--$13\times$ on a single nuclide at a
fixed energy) does not survive integration on a realistic
thermal pin cell. On Godiva ($3$ nuclides, fast spectrum)
the verdict is reversed --- SVD runs $1.22\times$ faster than
the table at the same $5.14\times$ memory penalty --- but
this is a small-problem win, not a general one. $k$-eigenvalue
fidelity is not the bottleneck: every provider lands inside
combined seed $\sigma$ of every other on both benchmarks, and
the engine reproduces the ICSBEP HMF-001 handbook value to
$+79 \pm 38$~pcm (inside the $\pm 100$~pcm experimental band).

\paragraph{Where SVD pays.}
The unambiguous SVD win in this engine is off-library
temperature operation. At a uniform $+150$~K offset on PWR,
SVD's three-point unity-normalised Ducru
reconstruction~\cite{ducru2017} produces $\sigma(E,T)$ from a
single basis at one library column; the pointwise table must
hold both bracketing library columns resident and
stochastically interpolate between them per
nuclide-collision. Memory inverts: SVD $156$~MB vs table
$206$~MB ($1.32\times$ smaller for SVD), CPU throughput
matches within $0.07\,\%$, and $k_\infty$ trails ACE$+$WMP by
$213$~pcm (about $3.6\sigma_\text{seed}$). The off-library
memory win is the one regime where SVD's reconstruction
discipline cleanly beats interpolation.

\paragraph{The representation-vs-engine memory gap.}
The literature reports SVD compression headlines as large as
$132.9\times$ vs.\ pointwise tables. That figure is correct as
a per-reaction, representation-only byte count. The
in-engine working set is a different quantity: it includes
overhead from discrete-level data (MT$=51$--$91$), URR
probability tables, S($\alpha,\beta$) inelastic distributions,
angular distributions, the geometry, and the BVH. Once those
are included, the SVD live-memory cost in our engine is
$5.12\times$ \emph{larger} than the pointwise baseline at
on-library $T$, not $132.9\times$ smaller. We document both
numbers and the reason they differ.

\paragraph{GPU is not the win either.}
On Godiva, GPU SVD on an RTX A1000 runs $1.3\times$ slower
than the same SVD code on a $20$-thread CPU rayon loop, and
the same factor slower than GPU pointwise on the same
benchmark. The kernel-level per-lookup advantage is
swamped by launch overhead and divergent memory access in
event-batched transport at this problem scale. GPU SVD
becomes competitive on photon-shielding workloads
($10^6$+ histories per case) and on per-nuclide
microbenches; it is not a Godiva or PWR k-eff win.

\paragraph{What we suggest reading from this.}
SVD compression of $\log_{10}\sigma(E,T)$ is interesting
research and is unambiguously useful in one operationally
common regime: off-library temperature operation in
multi-temp libraries. It is not, on the measurements
reported here, a general drop-in win against pointwise tables
or against ACE$+$WMP. The contribution of the paper is the
side-by-side measurement of the trade-offs, not a claim that
one provider should replace the others. The four-provider
interface remains live in the engine so future work --- on
problem classes we have not yet measured, on GPU hardware
we have not yet had access to, on multi-temp library setups
where the table baseline scales linearly in $T$ --- can run
its own A/B without re-doing the plumbing.
\end{abstract}

\vspace{0.3em}
\noindent\textbf{\small Keywords:} \small Monte Carlo,
cross-section compression, singular value decomposition,
windowed multipole, Doppler broadening, partition-of-unity
interpolation, nuclear data, GPU transport.
